\chapter{The Landscape}
\label{ch:landscape}

\begin{quote}
\itshape
Before you write a line of agent code, you need to know which of the four
Microsoft libraries with overlapping names you are supposed to be using, and
why the other three exist.
\end{quote}

\noindent
This chapter answers that question. It is the only chapter in the book with no
runnable sample in it, and it exists because the single most common way to waste
a week on Agent Framework is to start from a two-year-old blog post about a
library that has since been superseded.

By the end you should be able to state, without reference material, where Agent
Framework sits relative to \pkg{Microsoft.Extensions.AI}, Semantic Kernel,
AutoGen, and Microsoft Foundry --- and, more usefully, when you should not reach
for it at all.

% ============================================================
\section{Three years, compressed}
\index{Semantic Kernel}\index{AutoGen}

Microsoft arrived at Agent Framework by merging two projects that had been
solving adjacent halves of the same problem in different styles.

\textbf{Semantic Kernel} was the enterprise-facing effort. Its concerns were the
ones a platform team recognises immediately: dependency injection, plugins with
typed contracts, filters and telemetry, multi-language parity across .NET, Python
and Java, and a support story you could put in front of a risk committee. Its
weakness was that orchestration --- getting several agents to genuinely
collaborate --- was bolted on late.

\textbf{AutoGen} came from the research side, and orchestration was the whole
point. Group chat, conversational handoff, and agent topologies were first-class
there years before they were credible anywhere else. Its weakness was the mirror
image of Semantic Kernel's: excellent patterns, thin production plumbing.

Between them the two projects accumulated more than seventy-five thousand GitHub
stars and a large body of field experience. In October~2025 Microsoft announced
that they would converge into a single SDK. The timeline from there was short:

\begin{center}
\begin{tabular}{@{}ll@{}}
\toprule
\textbf{Date} & \textbf{Milestone} \\
\midrule
October 2025      & Agent Framework announced; convergence of SK and AutoGen \\
19 February 2026  & Release Candidate --- API surface frozen for 1.0 \\
2--3 April 2026   & \textbf{1.0 General Availability}, .NET and Python \\
June 2026 (BUILD) & Agent harness, hosted agents, CodeAct, handoff pattern \\
\bottomrule
\end{tabular}
\end{center}

The GA milestone is the one that matters for your purposes. It carries stable
APIs, versioned releases, and a long-term support commitment. Everything shipped
at BUILD in June builds on top of that foundation rather than changing it ---
which is why this book treats the GA surface as the spine and marks the June
additions separately.

\begin{note}
The predecessors have not vanished. Semantic Kernel's agent abstractions
continue to receive critical fixes for a period after GA, and AutoGen~.NET is in
maintenance --- it works, but it gets no new features. If you inherit a codebase
using either, you are not in an emergency. If you are starting something new,
there is no argument for either.
\end{note}

% ============================================================
\section{The layer cake}
\label{sec:layercake}
\index{layering}

The most useful mental model is a stack, and the most common confusion comes
from collapsing two of its layers into one.

\needdiagram{layer-cake}{The Agent Framework stack, from provider abstraction to hosting}{%
A four-band vertical stack. Bottom band: model providers (Azure OpenAI, OpenAI,
Anthropic, Bedrock, Gemini, Ollama). Second band: \texttt{Microsoft.Extensions.AI}
--- \texttt{IChatClient}, \texttt{IEmbeddingGenerator}. Third band: Semantic
Kernel foundations --- plugins, filters, kernel services. Fourth band: Microsoft
Agent Framework --- agents, tools, context providers, middleware, workflows. Top
band, drawn as an optional overlay to the side: Microsoft Foundry --- hosted
agents, memory store, evaluations. Arrows should run upward only.}

From the bottom:

\begin{description}[leftmargin=0pt, style=nextline]
\item[Model providers.] Azure OpenAI, OpenAI, Anthropic, Amazon Bedrock, Google
Gemini, and Ollama are all supported at 1.0. This layer is where your tokens are
actually spent.

\item[\pkg{Microsoft.Extensions.AI}.] The provider abstraction. Its central type
is \api{IChatClient}, and its job is to make the layer below interchangeable.
This is not part of Agent Framework --- it is a general .NET abstraction that
Agent Framework consumes. Understanding this separation is what makes the
one-line provider swap in Chapter~\ref{ch:first-agent} unsurprising rather than
magical.

\item[Semantic Kernel foundations.] Plugin infrastructure, filters, kernel
services. You will rarely program against this layer directly, but it is why
certain Agent Framework abstractions have the shape they do.

\item[Agent Framework.] Agents, tools, MCP integration, context providers,
middleware, and multi-step workflows. This is the subject of the book.

\item[Microsoft Foundry.] Not a layer so much as a managed environment that can
host what you build: hosted agents, a memory store, evaluation runs,
observability into Application Insights. Optional. Agent Framework runs perfectly
well with no Azure account at all, and most of this book does exactly that.
\end{description}

\begin{warning}
Agent Framework is \emph{built on} Semantic Kernel rather than being a
replacement for it, and people routinely get this backwards. If you read a claim
that Agent Framework ``replaces'' Semantic Kernel, the author means that it
replaces SK's \emph{agent abstractions} --- not the kernel, plugin and filter
infrastructure underneath.
\end{warning}

% ============================================================
\section{Two load-bearing abstractions}
\label{sec:abstractions}

Strip away the surface area and Agent Framework has two types you must
understand. Everything else is a convenience over one of them.

\subsection{\texttt{AIAgent}}
\index{AIAgent@\texttt{AIAgent}}

An agent is a model, a set of instructions, a set of tools, and a conversation.
The framework's expression of that is \api{AIAgent}, and the striking thing about
its construction is how little ceremony it requires: any \api{IChatClient}
becomes an agent through a single extension method.

That design choice has a consequence worth stating early. Because the agent is a
thin decoration over a chat client, an agent is cheap. You are not expected to
have one agent per application; you are expected to have several, each narrowly
scoped, composed by the second abstraction.

\subsection{\texttt{Workflow}}
\index{Workflow@\texttt{Workflow}}

A workflow is a directed graph of \emph{executors} connected by \emph{edges},
with typed messages passing along them. An agent is an executor. So is an
ordinary function. That symmetry is the framework's central idea, and it is what
separates Agent Framework from the majority of agent libraries, which model
composition as either a linear chain or a free-for-all conversation.

Because a workflow is a graph rather than a conversation, it can do things a
conversation cannot: fan out and rejoin, route conditionally, stream structured
events about its own progress, checkpoint mid-execution, and pause to wait for a
human decision before resuming hours later.

\begin{note}
If you take one thing from this chapter, take this: \emph{the orchestration
patterns are conveniences over the workflow graph}. Sequential, concurrent, group
chat and handoff (Chapter~\ref{ch:orchestration}) are not a separate mechanism.
Knowing that tells you what to do when a pattern stops fitting --- drop a level
and write the graph yourself.
\end{note}

% ============================================================
\section{What ships in the box}

The 1.0 feature set, stated plainly:

\begin{itemize}[leftmargin=1.4em]
\item \textbf{Multi-agent orchestration} --- sequential, concurrent, group chat
      and handoff patterns.
\item \textbf{Graph-based workflows} --- streaming, checkpointing and
      human-in-the-loop, available in both imperative and declarative form.
\item \textbf{Multi-provider support} --- six providers at GA, swappable
      through configuration.
\item \textbf{Extensible middleware} --- request and response interception with
      custom pipelines and exception handling.
\item \textbf{Built-in observability} --- OpenTelemetry integration following the
      GenAI semantic conventions, for distributed tracing and debugging.
\item \textbf{MCP} --- Model Context Protocol support for tools, native rather
      than bolted on.
\item \textbf{A2A} --- the Agent-to-Agent protocol, for structured messaging
      between agents across different runtimes. A Python agent and a .NET agent
      can coordinate.
\end{itemize}

The last two deserve emphasis because they arrived \emph{at} 1.0 rather than
after it. Cross-runtime and cross-vendor interoperability were treated as
release-blocking concerns, which is a reasonable signal about the framework's
intended lifespan.

% ============================================================
\section{The package map}
\label{sec:packages}
\index{NuGet packages}

Agent Framework is not one NuGet package. Knowing which is which saves a
surprising amount of time.

\begin{center}
\small
\begin{tabularx}{\linewidth}{@{}lX@{}}
\toprule
\textbf{Package} & \textbf{Purpose} \\
\midrule
\pkg{Microsoft.Agents.AI} & Core abstractions. \api{AIAgent}, sessions, tools, context providers, middleware. Always needed. \\
\pkg{Microsoft.Agents.AI.OpenAI} & OpenAI and Azure OpenAI provider binding. \\
\pkg{Microsoft.Agents.AI.Workflows} & Executors, edges, the workflow builder, checkpointing. Needed from Chapter~\ref{ch:workflows-basics} onward. \\
\pkg{Microsoft.Agents.AI.Workflows.Declarative} & YAML-defined workflows. \\
\pkg{Microsoft.Agents.AI.A2A} & Agent-to-Agent protocol client. \\
\pkg{Microsoft.Agents.AI.Foundry.Hosting} & Hosting an agent in Foundry Agent Service. \\
\bottomrule
\end{tabularx}
\end{center}

Additional provider packages (Anthropic, Bedrock, Gemini, Ollama), an ASP.NET
Core A2A hosting package, and a Roslyn source-generator package for compile-time
workflow route configuration are published alongside these. Appendix~\ref{app:packages}
carries the full list with the versions used throughout the book.

\begin{versionbox}
These packages ship in \emph{release trains}, not on one version line. As this
was written the core train --- \pkg{Microsoft.Agents.AI},
\pkg{Microsoft.Agents.AI.OpenAI}, \pkg{Microsoft.Agents.AI.Workflows} --- was at
\mafcore{}, while the extensions train carrying the harness, durable execution
and Foundry hosting was at \mafext{}, and two packages were still on a
\code{1.0.0-rc5} from March. A single \code{\$(MafVersion)} property will not
work. Pin every package explicitly; Central Package Management
(\code{Directory.Packages.props}) is set up in Chapter~\ref{ch:first-agent} for
exactly this reason, and Appendix~\ref{app:packages} lists every package with
its train.
\end{versionbox}

\needscreenshot{vs-nuget-browse-maf}%
{Browsing the Agent Framework packages in the NuGet Package Manager}%
{In \vsver{}, open \emph{Manage NuGet Packages} on any project, switch to the
\emph{Browse} tab, and search for \texttt{Microsoft.Agents.AI}. Capture the full
dialog including: the search box with the query visible, the result list showing
at least \texttt{Microsoft.Agents.AI}, \texttt{Microsoft.Agents.AI.OpenAI} and
\texttt{Microsoft.Agents.AI.Workflows}, the \emph{Include prerelease} checkbox,
and the right-hand version dropdown expanded if possible. Light theme preferred
for print legibility.}

% ============================================================
\section{A first look at the API}

You will not run anything until Chapter~\ref{ch:first-agent}. But the shape of
the API is worth seeing now, because it explains why the rest of the book spends
so little time on agent construction and so much on what agents are composed
into.

Here is an entire agent:

\begin{csharp}[caption={An agent, in full}, label={lst:first-look}]
using Microsoft.Agents.AI;

AIAgent agent = chatClient.AsAIAgent(
    instructions: "You are a helpful assistant.",
    name: "Assistant");

Console.WriteLine(await agent.RunAsync("What is Microsoft Agent Framework?"));
\end{csharp}

\begin{verifybox}
Listing~\ref{lst:first-look} is transcribed from Microsoft's published examples
rather than compiled by the author against \mafcore{}. The \api{chatClient}
variable is deliberately unbound here --- constructing a provider-specific
\api{IChatClient} is the first thing Chapter~\ref{ch:first-agent} does, and the
construction syntax differs by provider.
\end{verifybox}

Three observations, each of which the rest of the book expands on.

\textbf{The agent is an extension method away from a chat client.} There is no
agent builder, no registration ceremony, no configuration file. If constructing
an agent were expensive, composing a dozen of them into a workflow would be
impractical --- so it is not expensive.

\textbf{Instructions and name are constructor-level concerns.} The name is not
cosmetic. In handoff orchestration (Chapter~\ref{ch:orchestration}) the framework
generates transfer tools from agent names, so what you type there becomes part of
the model's decision surface.

\textbf{Running the agent returns a value you can print.} The synchronous,
non-streaming path is the simple case; streaming, sessions, tool approval and
cancellation all layer on top without changing this shape.

For contrast, here is the same idea at the other end of the book --- a three-agent
handoff topology, which is what Listing~\ref{lst:first-look} scales into:

\begin{csharp}[caption={Where this is going: a handoff topology (Chapter~\ref{ch:orchestration})}, label={lst:first-look-handoff}]
using Microsoft.Agents.AI;
using Microsoft.Agents.AI.Workflows;

AIAgent triage = chatClient.AsAIAgent(
    instructions: "You receive a user request and route it to the right specialist.",
    name: "Triage");

AIAgent billing = chatClient.AsAIAgent(
    instructions: "You handle billing questions.", name: "Billing");

AIAgent tech = chatClient.AsAIAgent(
    instructions: "You handle technical support questions.", name: "Tech");

Workflow workflow = AgentWorkflowBuilder
    .CreateHandoffBuilderWith(triage)
    .WithHandoff(triage, billing)
    .WithHandoff(triage, tech)
    .Build();
\end{csharp}

\begin{verifybox}
Also transcribed rather than compiled. \api{AgentWorkflowBuilder} lives in
\pkg{Microsoft.Agents.AI.Workflows}; verify the builder method names against
\mafworkflows{} before relying on them.
\end{verifybox}

Note what is \emph{not} in Listing~\ref{lst:first-look-handoff}: any routing
logic. You declare the topology and the guardrails; the agents make the routing
decisions at run time, using transfer tools the framework injects on your behalf.
That division --- topology to the developer, routing to the model --- is the
design position Agent Framework takes throughout, and Chapter~\ref{ch:orchestration}
examines where it holds up and where it does not.

% ============================================================
\section{When not to use Agent Framework}
\label{sec:when-not}

A book that never says no is a brochure. Four cases where reaching for Agent
Framework is the wrong call:

\begin{description}[leftmargin=0pt, style=nextline]
\item[You need one model call.] If your feature is ``summarise this text'' or
``classify this ticket'', you want \pkg{Microsoft.Extensions.AI} and an
\api{IChatClient} directly. An agent adds an instruction preamble, a tool-selection
step and a conversation abstraction that you are not using. The layer below is
right there; use it.

\item[Your control flow is deterministic.] If you know at design time exactly
which steps run in which order, that is a pipeline, and C\# already expresses
pipelines well. Workflows earn their complexity when routing depends on model
output, when steps can fail and resume, or when a human interrupts. If none of
those apply, you are paying for a graph engine to run a \code{foreach}.

\item[You are locked to a runtime Agent Framework does not target.] The packages
target .NET~8, .NET~Standard~2.0 and .NET Framework~4.7.2, which covers most
things --- but verify against your actual target before committing.

\item[Latency is your binding constraint and your task is tool-heavy.] Each
tool call in a traditional agent loop is a separate model turn. On a
tool-intensive task this dominates. Chapter~\ref{ch:harness} covers the CodeAct
approach to collapsing that loop --- available for .NET through
\pkg{Microsoft.Agents.AI.Hyperlight}, though still preview --- which materially
reduces the turn count without eliminating the model round trip. If you need a
chain of forty tool calls to return in under a second, no agent framework is your
answer yet.
\end{description}

\begin{note}
The inverse test is more useful than any of these. Agent Framework pays for
itself when at least two of the following are true: control flow depends on model
output; multiple specialised agents must collaborate; execution runs long enough
that a process restart is plausible; a human must approve or intervene mid-run;
or you need to explain afterwards, from traces, why the system did what it did.
\end{note}

% ============================================================
\section{Verifying your toolchain}
\label{sec:toolchain}

Chapter~\ref{ch:first-agent} builds the first project from nothing. Before
turning the page, confirm three things.

\subsection{SDK version}

\begin{shellcmd}
dotnet --list-sdks
\end{shellcmd}

You want \dotnetver{} or later on the machine you will develop on. The packages
themselves target lower, but the samples in this book use file-scoped namespaces,
collection expressions and, in two places, the \code{\#:package} file-based app
directive, which requires a recent SDK.

\subsection{IDE workloads}

Agent Framework needs no special Visual Studio workload beyond ASP.NET and web
development, plus the container tooling if you intend to follow
Chapters~\ref{ch:memory} and~\ref{ch:hosting}. Confirm your installation and
build number:

\needscreenshot{vs-about-dialog}%
{Visual Studio version used throughout this book}%
{\vsver{}: \emph{Help $\rightarrow$ About Microsoft Visual Studio}. Capture the
dialog showing the full version and build number, and the edition. This appears
once, in this chapter, so that readers on a different build know exactly what the
later screenshots were taken against. Crop tightly to the dialog.}

\subsection{Solution layout}

Every sample in this book lives in one solution, organised by chapter. Setting it
up is the first task of Chapter~\ref{ch:first-agent}, but the finished shape is
worth seeing now so you know what you are building toward:

\needscreenshot{vs-solution-explorer-scaffold}%
{The companion solution, organised by chapter}%
{\vsver{} Solution Explorer, showing the \texttt{maf-book-samples} solution
expanded two levels: solution folders \texttt{01-agents}, \texttt{02-tools},
\texttt{03-memory}, \texttt{04-workflows}, \texttt{05-orchestration},
\texttt{06-observability}, \texttt{07-hosting}, \texttt{08-harness}, with at
least one project visible inside \texttt{01-agents}, and
\texttt{Directory.Packages.props} and \texttt{Directory.Build.props} visible at
solution level. Collapse everything below project level so the structure is
legible at print size.}

% ============================================================
\section{Summary}

Agent Framework is the converged successor to Semantic Kernel's agent
abstractions and AutoGen's orchestration patterns, GA since April~2026, with
stable APIs and long-term support. It sits above
\pkg{Microsoft.Extensions.AI} --- which supplies the interchangeable
\api{IChatClient} --- and below Microsoft Foundry, which can host what you build
but is not required to run it.

Two abstractions carry the framework. \api{AIAgent} is a deliberately cheap
decoration over a chat client. \api{Workflow} is a typed, checkpointable graph of
executors, and every orchestration pattern in the framework is a convenience over
it.

It is the wrong tool for a single model call, for deterministic pipelines, and
for latency-bound tool chains. It is the right tool when control flow depends on
model output, when several agents collaborate, when runs are long enough to need
resumption, and when you must be able to explain afterwards what happened.

\begin{exercisebox}
Answer these from memory before continuing. If you cannot, reread
\S\ref{sec:layercake} and \S\ref{sec:abstractions} rather than pressing on ---
Chapter~\ref{ch:first-agent} assumes all four.

\begin{enumerate}[leftmargin=1.4em]
\item In three sentences, where does Agent Framework sit relative to
      \pkg{Microsoft.Extensions.AI}, Semantic Kernel, and Microsoft Foundry?
\item Why is it significant that an agent is constructed from an
      \api{IChatClient} by a single extension method?
\item What is the relationship between the four orchestration patterns and the
      workflow graph, and what does that relationship tell you to do when a
      pattern does not fit your problem?
\item Give two concrete features of your own current or previous work where
      Agent Framework would have been the wrong choice, and say which layer you
      would have used instead.
\end{enumerate}
\end{exercisebox}
